\section{Exhausting the endoscopic spectrum in every even degree}
\label{sec:kt-endoscopic-converse}

This section extends KT-U3-ENDO-INJECT to an equality. Its additional
published inputs are the rank-three spectral classification of Rogawski
and the spherical representation theorem for a special maximal parahoric.
The latter is used to prove the converse local-level assertion, rather
than to assume that a compact fixed vector automatically gives an
unramified base change.

\begin{lemma}[KT-U3-LEVEL-CONVERSE]
Every irreducible smooth representation of $G(\mathbf Q_3)$ with a
vector fixed by the full standard lattice stabilizer has base change to
$\mathrm{GL}_3(E_3)$ with unramified Weil--Deligne parameter and zero
monodromy. For a fixed such local parameter, at most one member of the
packet has a full-stabilizer fixed vector, and its fixed space is one
dimensional.
\end{lemma}
\begin{proof}
Let $K_3^\circ$ be the connected special maximal parahoric inside $K_3$.
Adler--Lansky Sections~4.1--4.2.1, pp.~10--12, identify their reductive
quotients with $\mathrm{SO}_3$ and $O_3$, respectively.
Haines, \emph{On Satake parameters for representations with parahoric
fixed vectors}, author version~4, Section~3, Proposition~3.1, p.~4,
shows that every $K_3^\circ$-spherical irreducible representation is the
spherical constituent of a principal series induced from a weakly
unramified character of the minimal Levi. This is the ordinary
representation-theoretic assertion of that proposition; no conjectural
compatibility between Haines's Satake parameter and local Langlands is
being used.

Let $I(\lambda)$ be such a principal series containing our representation.
Taking invariants under the compact open subgroup $K_3$ is exact in
characteristic zero, so $I(\lambda)^{K_3}\ne0$. The Iwasawa decomposition
$G(\mathbf Q_3)=B(\mathbf Q_3)K_3$ identifies this space with the
functions on one coset satisfying the inducing character relation.
It is therefore one dimensional, and forces $\lambda$ to be trivial on
$B\cap K_3$. In the coordinates of Adler--Lansky (4.4.1), this says
that $\lambda_1$ is unramified on $E_3^\times$ and $\lambda_2=1$ on
$E_3^1$.

Adler--Lansky Section~4.4 now leaves only the irreducible principal
series, case (4.4PS--1), and case (4.4PS--3). Case (4.4PS--2) is
impossible because its restriction to $\mathbf Q_3^\times$ has ramified
quadratic part, whereas the restriction of $\lambda_1$ is unramified.
In (4.4PS--1), the constituents are the trivial representation and the
Steinberg representation. The trivial constituent already has a fixed
vector; exactness and $\dim I(\lambda)^{K_3}=1$ show that the Steinberg
constituent has none. In (4.4PS--3), exactly one of the two constituents
has a fixed vector, again by exactness.

The full list of base changes is Proposition~7.1(i)--(iv), p.~28.
In case (i), it is the irreducible unramified principal series of
$\mathrm{GL}_3(E_3)$. In case (ii), where the unitary-group principal
series is irreducible but the general-linear principal series is
reducible, the displayed lift is induced from a one-dimensional
$\mathrm{GL}_2$ representation and an unramified character; its
parameter is again a direct sum of unramified characters with zero
monodromy. In case (iii), the full-stabilizer constituent is the trivial
representation and its lift is trivial, also with zero monodromy.
Case (iv) is the irreducible unramified general-linear principal series.
These cover every case above. They are the explicit principal-series
cases, whose proof refers to the earlier unramified Proposition~4.1
(Proposition~5.1 in that author preprint); neither the depth-preservation
hypothesis nor restrictions on the size of the residue field are used.

For a fixed base-change parameter the unordered unramified inducing
triple determines $\lambda_1$ up to inversion, hence determines
$I(\lambda)$ up to the Weyl action. The same fixed-space calculation
then shows that at most one constituent in its packet is $K_3$-spherical,
with a one-dimensional fixed space. This includes the reducible
quadratic case, where the two members share one principal series.
\end{proof}

\begin{theorem}[KT-U3-ENDO-EXACT]
For every even $e\ge2$, the full compact space $M_{e-1}$ is a direct
sum of a stable part and the endoscopic lines constructed in
KT-U3-ENDO-INJECT. Its endoscopic dimension is exactly
\[
             E_e=\left\lfloor\frac{e-1}{3}\right\rfloor,
\]
and its stable dimension is
\[
 S_e=D(e-1)-E_e=
 \begin{cases}
 e^3/216-e/3+1,&e\equiv0\pmod6,\\
 e^3/216-2e/9+11/27,&e\equiv2\pmod6,\\
 e^3/216-2e/9+16/27,&e\equiv4\pmod6.
 \end{cases}                                                   \tag{KT-E.3}
\]
Here stable means that the generic parameter is a single cuspidal
$\mathrm{GL}_3/E$ constituent. Each endoscopic parameter contributes
exactly one line, and is uniquely of the form
$1\boxplus\operatorname{BC}_{E/\mathbf Q}(\pi_f)$ for a non-CM
quadratic-twist orbit in $S_{2e+1}(\Gamma_0(3),\epsilon)$.
\end{theorem}
\begin{proof}
\paragraph{The rank-three spectral input and regularity.}
We use Rogawski's classification for unitary groups in three variables,
including their inner forms: the discrete spectrum is partitioned by
discrete global Arthur parameters, with the expected infinitesimal
characters and local base changes. This rank-three assertion is a
proved published input, not the unrestricted general-rank Arthur
formalism. An explicit source locator is Bella\"iche--Chenevier,
\emph{Families of Galois representations and Selmer groups},
Appendix~A.9, p.~244 in the final author version
(pp.~228--229 in the supplied preprint): the paragraph
preceding Remark~A.9.1 states that this spectral description was proved
by Rogawski for $m\le3$, referring to his monograph, Chapters~13--14.
Remark~A.9.1 records the compact regularity obstruction. Section~6.7,
p.~148 (p.~140 in the supplied preprint), defines regularity as strict inequalities between coefficient
highest weights.

In our case the coefficient weights are $(e-1,0,1-e)$, strictly
decreasing for $e\ge2$. One can check the required obstruction directly
in rank three. A nontrivial Arthur $\mathrm{SL}_2$ block must have size
two or three. A size-two block places a pair of weights differing by
one in the infinitesimal character; a size-three block places three
consecutive weights there. The required multiset $(e,0,-e)$ has no
such pair. Therefore every parameter contributing to this space is
generic. The only remaining partitions of the rank are $3$, $2+1$,
and $1+1+1$. For these generic parameters we use the already checked
KMSW Theorem~5.0.5, including its multiplicity formula for this pure
inner form.

At each finite place, the standard compact subgroup forces the
base-change parameter to be unramified with zero monodromy:
this is the usual spherical local correspondence at split and
unramified nonsplit places, and KT-U3-LEVEL-CONVERSE at three.
Consequently every simple global constituent is unramified at every
finite place of $E$. The determinant of the rank-three parameter is
an everywhere unramified Hecke character with trivial infinity type.
Since $E$ has class number one, this determinant is trivial.

\paragraph{All possible one-dimensional constituents.}
For an integer $s$, an everywhere unramified Hecke character of $E$
with infinity character $(z/\bar z)^s$ exists exactly when $3\mid s$,
and is then unique; denote it by $\theta_s$.
Indeed, the class number is one, so its values on ideals are forced
by a generator. The rule is well defined on different generators
exactly when $(u/\bar u)^s=1$ for every $u\in\mu_6$, which is the
condition $3\mid s$. This also proves uniqueness; $\theta_0=1$.
Changing the reciprocity convention inverts all these characters
and leaves the assertions below unchanged.

At an inert rational prime the local value on a uniformizer is one,
because that uniformizer can be chosen rational. At the ramified
prime three, use a generator proportional to $\sqrt{-3}$; the ratio
of it to its conjugate is $-1$, so
\[
                  \theta_s(\varpi_3)=(-1)^s.
\]
In particular $\theta_e$ and $\theta_{-e}$ are trivial at every
finite nonsplit place whenever they exist, since $e$ is even.

\paragraph{The middle character gives precisely the classical lifts.}
Suppose the generic parameter has partition $2+1$ and its
one-dimensional summand occupies the middle infinity weight zero.
That summand is $\theta_0=1$, so write the parameter as $1\boxplus\Pi$.
The cuspidal $\mathrm{GL}_2/E$ representation $\Pi$ has determinant
one. The rank-two duality identity therefore gives
$\Pi^\vee\simeq\Pi$. Since it is conjugate self-dual, also
$\Pi^c\simeq\Pi$. Langlands, \emph{Base change for $\mathrm{GL}(2)$},
Section~2(B),(C),(E), p.~13, and Lemma~11.6(b), p.~142,
then give exactly two cuspidal descents to $\mathrm{GL}_2/\mathbf Q$,
related by twisting with $\epsilon$.

Their real parameter restricts to $(z/\bar z)^{\pm e}$, so their
classical weight is $2e+1$. Their central characters have trivial
base change, hence are either $1$ or $\epsilon$. The odd weight
forces the odd character $\epsilon$.
At every finite prime other than three, local base change is
unramified only if the original inertia and monodromy are trivial:
at a nonsplit unramified prime the inertia groups are identical,
and at a split prime the assertion is immediate.
At three the restriction to $W_{E_3}$ is unramified with zero
monodromy, so inertia factors through the quadratic quotient.
Its determinant is $\epsilon$, hence its characters are exactly
$1$ and $\epsilon$. Thus the conductor is one at three and zero
elsewhere. The descents are precisely the level-three primitive
forms in KT-U3-ENDO-INJECT. They are non-CM because their base
change is cuspidal.

\paragraph{Outer characters and three-character parameters are excluded.}
If a $2+1$ parameter has its one-dimensional constituent at an
outer weight, that constituent is $\theta_e$ or $\theta_{-e}$,
so $3\mid e$. Let $a$ act by $-I_2$ on its cuspidal rank-two
constituent and by $1$ on the character. At every finite nonsplit
place the character is trivial, and the determinant of the
rank-two constituent is one. Its unramified roots are therefore
$b,b^{-1}$. Exactly the connected determinant-one centralizer
argument in KT-U3-ENDO-INJECT puts $a$ in
$\operatorname{Cent}(\psi_v,\mathrm{SL}_3)^\circ$.
This argument does not require a bound on $|b|$.
At split places the same element belongs to the connected
determinant-one centralizer of the semisimple parameter.
Every finite local packet character consequently has value one
on $a$.

At infinity the rank-two constituent now occupies adjacent
positions, either $(1,2)$ or $(2,3)$, in the ordered weights
$(e,0,-e)$. The compact-root rule from Gan--Gross--Prasad,
Section~2, p.~115, gives value $(-1)^{1+2}=-1$ or
$(-1)^{2+3}=-1$. The global product is nontrivial, whereas
$\epsilon_\psi=1$ for a generic parameter. The multiplicity
formula excludes this tensor product.

For partition $1+1+1$, the only possible parameter is
$\theta_e\boxplus1\boxplus\theta_{-e}$, again requiring
$3\mid e$. At every finite nonsplit place all three characters
are trivial. Choose the global sign element acting as $-1$ on
two adjacent infinity positions and as $1$ on the third.
It lies in $\mathrm{SO}_3(\mathbf C)$, hence in the connected
determinant-one centralizer, at every such finite place. At
split places it belongs to the connected determinant-one torus
centralizer. All finite character values are one, and the
compact real value is minus one. The same generic multiplicity
formula excludes this parameter from the compact space.

\paragraph{Multiplicity and dimensions.}
The preceding cases exhaust the rank-three classification.
Every surviving endoscopic parameter is therefore one of the
middle-character parameters already constructed. At the real
place it has only one member on the compact group with this
coefficient type. At finite split and unramified nonsplit places
the spherical member and its fixed line are unique; at three
this is KT-U3-LEVEL-CONVERSE. KMSW gives multiplicity one for
the resulting global tensor product. Thus each non-CM twist
orbit contributes exactly one line and no further endoscopic
lines occur. KT-U3-ENDO-INJECT gives their number as
$\lfloor(e-1)/3\rfloor$. Subtracting this from KT-U3-DIM in
the three even residue classes modulo six gives (KT-E.3).
\end{proof}

\begin{corollary}[KT-U3-SPECTRAL-SPLIT]
Let $p\ne3$ be split in $E$, and write $Q_{e,p}(Y)$ for the
characteristic polynomial of the classical $T_p$ on
$S_{2e+1}(\Gamma_0(3),\epsilon)$. Put
\[
 C_{e,p}(Y)=
 \begin{cases}
 1,&3\nmid e,\\
 Y-(\varpi^{2e}+\bar\varpi^{2e}),&3\mid e,
 \end{cases}
 \qquad \varpi\bar\varpi=p.
\]
The endoscopic polynomial $P^{\rm end}_{e,p}(X)$ is the unique
monic polynomial satisfying
\[
 \bigl(P^{\rm end}_{e,p}(X)\bigr)^2
       =\frac{Q_{e,p}(X-p^e)}{C_{e,p}(X-p^e)}.             \tag{KT-E.4}
\]
It belongs to $\mathbf Z[X]$. If $P^{\rm full}_{e,p}$ is the
characteristic polynomial of the scaled compact operator
$p^{e-1}T_p$, then
\[
 P^{\rm st}_{e,p}=P^{\rm full}_{e,p}/P^{\rm end}_{e,p}
       \in\mathbf Z[X]                                  \tag{KT-E.5}
\]
is the polynomial on the stable part. In particular the exact
uniform character algorithm at $p=7$, together with the classical
algorithm, computes the two factors separately in every even degree.
\end{corollary}
\begin{proof}
KT-GL2-CM gives exactly the displayed CM factor. When it occurs,
$2e$ is divisible by six, so its value is independent of the unit
chosen in a generator $\varpi$ of the split prime ideal.
Every other normalized classical newform occurs in a two-element
quadratic-twist orbit. At a split prime $\epsilon(p)=1$, so the
two eigenvalues $a_p$ coincide. Removing the CM factor therefore
leaves the square of a monic polynomial, and KT-E.1 shifts each
orbit eigenvalue by $p^e$. The quotient in (KT-E.4) is in
$\mathbf Z[X]$; comparing coefficients in its monic square root
puts that root in $\mathbf Q[X]$, while its roots are algebraic
integers. Hence the square root lies in $\mathbf Z[X]$.
KT-U3-ENDO-EXACT identifies it with the entire endoscopic factor.
The full polynomial is integral by KT-U3-SPECTRUM (or the same
scaled-lattice argument at the chosen split prime). Polynomial
division by the monic endoscopic factor gives the stable quotient
in $\mathbf Q[X]$; its integral roots again put it in
$\mathbf Z[X]$.
\end{proof}

\paragraph{Refinement of rank-three source coverage.}
An admissible rank-three source in K17-U3-GENERAL-COVERAGE
is attached to a cuspidal $\mathrm{GL}_3/E$ representation.
Its compact transfer therefore lies in the stable summand.
Consequently KT-TRACE-MULTIPLICITY remains valid with the
rank-three factor $P^{\rm full}_{e,7}(7^{2e}+7^e+1)$ replaced
by $P^{\rm st}_{e,7}(7^{2e}+7^e+1)$. The classical factor
is retained unchanged. This removes the identified endoscopic
factors from that comparison but makes no assertion that the
remaining stable comparison values have only small prime divisors.

\paragraph{Scope of the new source checks.}
The complete arguments above were checked in this run. The
published rank-three classification and the local/global
Langlands correspondences remain imported theorems; no claim
is made to have reconstructed Rogawski's trace formula or
compiled a new formal proof of it. The generic multiplicity
formula, its pure-inner scope, and the exact ramified
principal-series cases were checked at their printed
statements and used proofs as recorded in the check log.
